% Generated draft: requires mathematical and source review.
\documentclass[11pt]{article}
\usepackage[margin=25mm]{geometry}
\usepackage{fontspec}
\setmainfont{Latin Modern Roman}
\usepackage{amsmath,amssymb,mathtools,unicode-math}
\setmathfont{Latin Modern Math}
\usepackage{enumitem,needspace,xurl,hyperref}
\setlength{\parindent}{0pt}
\setlength{\parskip}{5pt}
\setlength{\emergencystretch}{3em}
\newcommand{\N}{\mathbb{N}}
\newcommand{\Z}{\mathbb{Z}}
\newcommand{\Q}{\mathbb{Q}}
\newcommand{\R}{\mathbb{R}}
\newcommand{\C}{\mathbb{C}}
\newcommand{\F}{\mathbb{F}}
\newcommand{\A}{\mathbb{A}}
\newcommand{\PP}{\mathbb{P}}
\newcommand{\E}{\mathbb{E}}
\newcommand{\eps}{\varepsilon}
\newcommand{\dd}{\,\mathrm{d}}
\begin{document}
\section*{UnsolvedMath ID: 11000223}
\textbf{Source problem code:} catalogue entry
\subsection{Definitions and mathematical statement}
In the usual closed oriented genus-\(g\) surface setting, let \(H=H^1(S;\R)\), with symplectic form \(\Omega(u,v)=\langle u\smile v,[S]\rangle\). A subspace \(L\subseteq H\) is Lagrangian if \(\dim L=g\) and \(\Omega|_{L\times L}=0\). The flattened \(H^1\) typography and exact surface convention must be checked; the dual homology formulation requires its own identification.

\emph{Problem 16 (dynamical Hodge decomposition).} Study the distributions of Lagrangian subspaces in \(H^1(S;\R)\) defined by Teichm\"uller geodesic flow, particularly their continuity. Is there a topological or geometric construction of these subspaces? Here ``distribution'' is naturally read as a field of subspaces over the relevant flow parameter space, but the base space and actual dynamical definition are absent.

If the intended construction assigns \(L_x\) over points \(x\) of a specified base \(B\), continuity should be expressed as continuity of the corresponding section of its Lagrangian Grassmannian bundle, after the appropriate local cohomology identifications. The excerpt does not say whether the field is defined at every point, only almost everywhere, or on a selected invariant subset, nor which topology is intended. The summary suggests the Kontsevich--Zorich cocycle, but does not identify which invariant subspaces are meant, specify an invariant measure, or supply nonzero-exponent hypotheses ensuring Lagrangian dimension. A measurable almost-everywhere dynamical splitting is not silently promoted to a continuous globally defined one. The source asks for intrinsic topological or geometric meaning as well as regularity, without providing a proposed construction.
\subsection{Short English statement}
Investigate the flow-defined Lagrangian subspace fields in surface cohomology, especially their continuity and whether they have a topological or geometric definition. The precise dynamical construction, base, and domain of definition are omitted.
\subsection{Sources}
\begin{itemize}
\item Local UnsolvedMath registry, permanent ID 11000223, source code catalogue entry, Problem 16, ``Dynamical Hodge decomposition,'' PDF page 262. No individual proposer is supplied.
\item Farb, \emph{Problems on Mapping Class Groups and Related Topics}. Supplied URL: \url{https://www.math.uchicago.edu/~farb/papers/mcgbook.pdf}; recorded access 2026-07-29; extraction pdf-font-structure. The edited volume contains 329 explicitly typeset headings across chapters, with repeated numbering disambiguated by page.
\item The supplied summary mentions a relationship with the Kontsevich--Zorich cocycle but says it cannot verify the exact problem. No further authoritative definition or full reference is supplied.
\item Supplied status open and OPEN-TRIAGE. Current primary-literature status and redistribution terms for the publicly accessible author/publisher copy remain NEEDS\_REVIEW.
\end{itemize}
\end{document}
